\documentclass[12pt,a4paper]{article}

\usepackage[utf8]{inputenc}
\usepackage[T1]{fontenc}
\usepackage{lmodern}
\usepackage[margin=2.3cm]{geometry}
\usepackage{xcolor}
\usepackage{setspace}
\usepackage{enumitem}
\usepackage{tcolorbox}
\usepackage{amsmath}
\usepackage{fancyhdr}
\usepackage{titlesec}
\usepackage[hidelinks]{hyperref}

\definecolor{PrimaryColor}{HTML}{213448}
\definecolor{SecondaryColor}{HTML}{547792}
\definecolor{AccentColor}{HTML}{649CA5}
\definecolor{RayColor}{HTML}{E69650}
\definecolor{DarkGray}{HTML}{555555}

\setstretch{1.25}
\setlength{\parskip}{0.55em}
\setlength{\parindent}{0pt}

\pagestyle{fancy}
\fancyhf{}
\fancyhead[L]{\small\color{DarkGray}Fractals --- Speaker Scripts}
\fancyhead[R]{\small\color{DarkGray}CSE 200}
\fancyfoot[C]{\small\color{DarkGray}\thepage}
\renewcommand{\headrulewidth}{0.4pt}

\titleformat{\section}{\Large\bfseries\color{PrimaryColor}}{}{0pt}{}
\titleformat{\subsection}{\large\bfseries\color{SecondaryColor}}{}{0pt}{}

% click cue -- tells the speaker to advance the slide
\newcommand{\cue}{\textcolor{RayColor}{\textsf{\footnotesize[CLICK]}}\hspace{0.3em}}
% a word to lean on while speaking
\newcommand{\stress}[1]{\textbf{#1}}

\newtcolorbox{speakerbox}[1]{
  colback=PrimaryColor!6, colframe=PrimaryColor, arc=2mm,
  boxrule=1pt, left=8pt, right=8pt, top=6pt, bottom=6pt,
  title={\bfseries #1}, fonttitle=\normalsize}

\newtcolorbox{notebox}[1]{
  colback=AccentColor!8, colframe=AccentColor, arc=2mm,
  boxrule=0.8pt, left=8pt, right=8pt, top=5pt, bottom=5pt,
  title={\bfseries #1}, fonttitle=\small}

\begin{document}

% ================================================================
\begin{center}
  {\Huge\bfseries\color{PrimaryColor} Fractals}\\[6pt]
  {\large\color{SecondaryColor} Infinite Detail from Simple Rules}\\[14pt]
  {\Large\bfseries Speaker Scripts}\\[6pt]
  {\color{DarkGray} Three speakers $\cdot$ three minutes each $\cdot$ CSE 200 Presentation}
\end{center}

\vspace{0.6cm}

\section*{How to use this}

Each script below is written to be \stress{spoken, not read silently}. The
sentences are short on purpose --- short sentences are much easier to deliver
when you are nervous.

\begin{itemize}[leftmargin=1.2em, itemsep=3pt]
  \item \cue means \stress{advance the slide}. Press the key, \emph{then} keep talking.
  \item \stress{Bold words} are the ones to lean on. Slow down slightly when you say them.
  \item Every script is about 390--410 words. At a calm speaking pace
        (\raisebox{0pt}{$\sim$}130 words a minute) that is just about 3 minutes.
        \stress{Time yourself once out loud} --- reading in your head is far faster
        and will fool you.
  \item You do not have to memorise this word for word. Learn the \stress{order of
        the ideas}. If you say them in your own words, that is better, not worse.
\end{itemize}

\begin{notebox}{The single most useful tip}
  When you finish a sentence, \stress{stop}. Do not fill the gap with ``um'' or
  ``so''. A one-second silence feels enormous to you and completely normal to the
  audience. It is the fastest way to sound like you know what you are doing.
\end{notebox}

\subsection*{Timing at a glance}

\begin{center}
\begin{tabular}{@{}p{2.4cm}p{5.6cm}p{2.0cm}p{4.4cm}@{}}
\hline
\textbf{Speaker} & \textbf{Covers} & \textbf{Time} & \textbf{Your one job} \\
\hline
Speaker 1 \newline {\footnotesize Aditya} & History, definition, Sierpinski triangle & 3 min &
  Show that normal measurement \emph{breaks}. \\[4pt]
Speaker 2 \newline {\footnotesize Sami} & Fractal dimension & 3 min &
  Give the broken thing a \emph{number}. \\[4pt]
Speaker 3 \newline {\footnotesize Pritom} & Nature, trees, the body, ferns, game worlds & 3 min &
  Show where that number \emph{shows up}. \\
\hline
\end{tabular}
\end{center}

\subsection*{Names you will have to say out loud}

\begin{center}
\begin{tabular}{@{}p{5.2cm}p{4.6cm}@{}}
\hline
\textbf{Name} & \textbf{Say it as} \\
\hline
Karl Weierstrass      & VY-er-shtrass \\
Georg Cantor          & GAY-org CAN-tor \\
Helge von Koch        & HEL-ga fon KOKH \\
Wac\l{}aw Sierpi\'nski & VAT-swaf sher-PIN-skee \\
Felix Hausdorff       & HOWS-dorf \\
Gaston Julia          & gas-TOHN ZHOO-lya \\
Pierre Fatou          & pee-AIR fa-TOO \\
Benoit Mandelbrot     & ben-WAH MAN-del-brot \\
Michael Barnsley      & BARNZ-lee \\
Albert Bosman         & BOSS-man \\
Romanesco             & roh-ma-NESS-koh \\
\hline
\end{tabular}
\end{center}

\begin{notebox}{If you stumble on a name}
  Say it confidently and move on. Nobody in the room will correct you, and
  stopping to fix it costs you five seconds and your rhythm.
\end{notebox}

\newpage
% ================================================================
\section*{Speaker 1 (Aditya) --- The Surprise}

\begin{speakerbox}{Your goal in one sentence}
  Convince the audience that ordinary geometry \stress{completely breaks} on
  these shapes --- so that Speaker 2 has a problem worth solving.\\[4pt]
  {\small\color{DarkGray} Slides: A Crisis in 19th-Century Mathematics $\rightarrow$
  Handoff to Speaker 2 \quad$\cdot$\quad about 400 words $\approx$ 3 min}
\end{speakerbox}

\vspace{0.3cm}

Think about zooming in on a curve. Any curve --- a circle, a parabola. Zoom in
far enough and it straightens out --- it looks like a straight line. For about two
hundred years, all of calculus quietly assumed this is \stress{always} true.

\cue In 1872, Karl Weierstrass broke that assumption --- on purpose. He built a
function that is \stress{continuous everywhere}, but \stress{differentiable
nowhere}. It is infinitely jagged --- zoom in as far as you like, it never
straightens out. No tangent line anywhere.

\cue Mathematicians called these shapes \stress{monsters}. But once the door was
open, more came. Cantor built a set of infinitely many points with zero total
length. Von Koch drew a curve that is nowhere smooth. Sierpi\'nski made shapes
built out of smaller copies of themselves. And Hausdorff gave these shapes a
\stress{number} --- a fractional dimension, in between the whole numbers we know.

For almost a hundred years, all of this was considered \stress{beautiful but
useless}.

\cue Then, in the 1960s, Benoit Mandelbrot --- working at IBM --- asked a very
simple question. \stress{How long is the coast of Britain?}

The honest answer is: it depends on your ruler. Measure with a hundred-kilometre
ruler, you get one number. Measure with a one-kilometre ruler, and the coast gets
\stress{longer} --- every little bay now counts. The length never settles down.

Mandelbrot's insight was that this is not a flaw in the coastline. It means
\stress{length is the wrong tool}. And the right tool was exactly that ``useless''
mathematics. He also had something Cantor and Koch never had --- \stress{computers}
--- so he could finally \emph{see} these shapes. He gave them one name:
\stress{fractals}, from the Latin \emph{fractus}, meaning ``broken''.

\cue So what is a fractal? Formally: a shape with detailed structure at
arbitrarily small scales, whose fractal dimension is greater than its topological
dimension. In plain words --- \stress{it has detail at every zoom level, and it is
built by repeating one simple rule}.

\cue Here is the classic example. Take a triangle. Remove the middle one.
Repeat --- forever. That is the Sierpi\'nski triangle.

\cue Now watch the numbers. Every step keeps three quarters of the area. So after
$n$ steps the area is $(3/4)^n$ --- and that goes to \stress{zero}.

\cue But the perimeter gets multiplied by three-halves at every step. That goes to
\stress{infinity}.

\cue So we have a two-dimensional shape with \stress{zero area} and
\stress{infinite perimeter}. Our normal tools have completely broken down.

Over to Speaker 2 --- if area and perimeter are useless here, what number
\emph{do} we use?

\begin{notebox}{Notes for delivery}
  \begin{itemize}[leftmargin=1.2em, itemsep=2pt, topsep=2pt]
    \item You have the most slides of the three. \stress{Do not linger} on the
      list of names --- say them as one quick sweep and move on.
    \item The area and perimeter slides have the maths already written. Do not read
      the algebra out loud. Just say the result: shrinks to zero, grows to infinity.
    \item \stress{If you blank out}, this one sentence carries your whole section:
      ``There are shapes with zero area but infinite perimeter, and normal geometry
      cannot handle them.''
  \end{itemize}
\end{notebox}

\newpage
% ================================================================
\section*{Speaker 2 (Sami) --- The Number}

\begin{speakerbox}{Your goal in one sentence}
  Take the thing Speaker 1 broke, and \stress{measure it} --- turn ``rough'' from a
  word into a number.\\[4pt]
  {\small\color{DarkGray} Slides: Fractal Dimension section $\rightarrow$
  Handoff to Speaker 3 \quad$\cdot$\quad about 390 words $\approx$ 3 min}
\end{speakerbox}

\vspace{0.3cm}

So we are stuck. Area says zero. Perimeter says infinity. Neither one tells us
anything useful. We need a different kind of measurement --- and the way in is to
rethink what \stress{dimension} actually means.

\cue Forget the usual definition for a moment. Let us define dimension by what
happens when you \stress{scale} something.

Take a line, and shrink it by half. How many of those copies fit in the original?
\stress{Two}.

\cue Take a square, and shrink each side by half. \stress{Four} copies fit.

\cue Take a cube, shrink it by half. \stress{Eight} copies.

Two, four, eight. That is $2^1$, $2^2$, $2^3$. Look at the exponent --- one, two,
three. \stress{The exponent is the dimension.}

\cue So in general: if we scale by a factor $r$, the number of self-similar pieces
is $N = (1/r)^D$. And if we solve that for $D$, we get
$D = \log N / \log(1/r)$.

Let us sanity-check it. A line: $\log 2 / \log 2 = 1$. A square:
$\log 4 / \log 2 = 2$. A cube: $\log 8 / \log 2 = 3$. It agrees with everything we
already believe. Good --- the formula is not doing anything strange.

\cue Now let us point it at the Sierpi\'nski triangle. Shrink it by half. How many
copies do we get? Not four --- \stress{three}. The middle one is missing.

So $D = \log 3 / \log 2$, which is about \stress{1.585}.

That is not a whole number. The Sierpi\'nski triangle is \stress{more than a line,
but less than a solid surface}. And that is exactly why area and perimeter failed
--- we were measuring a 1.585-dimensional object with one-dimensional and
two-dimensional tools.

\cue The same method works on the Koch curve. Scale it by one third, and you get
four copies. So $D = \log 4 / \log 3$, about \stress{1.262}.

\cue There is one catch. That formula needs \stress{perfect} self-similarity, and
a real coastline is not perfectly self-similar. So we use the general version: the
\stress{box-counting dimension}. Cover the shape with a grid of boxes of size
$\varepsilon$, count how many boxes contain part of the shape, then let
$\varepsilon$ shrink. The dimension is the limit of
$\log N / \log(1/\varepsilon)$.

Same idea --- but this one works on anything. Even a photograph.

\cue So here is what we have gained. \stress{Roughness is no longer just a word.
It is a number.} We can now say that one coastline is rougher than another, and
say \emph{by how much}.

Over to Speaker 3 --- where does this number actually show up?

\begin{notebox}{Notes for delivery}
  \begin{itemize}[leftmargin=1.2em, itemsep=2pt, topsep=2pt]
    \item The line-square-cube build-up is the heart of your section. \stress{Go
      slowly there.} If they follow two, four, eight, they will follow everything
      after it.
    \item Say ``log three over log two'', not ``logarithm of three divided by''.
      Shorter is clearer out loud.
    \item You do not need to explain what a logarithm is. Nobody will ask.
    \item \stress{If you blank out}: ``Dimension is just the exponent that tells you
      how many copies appear when you shrink a shape --- and for fractals that
      exponent is not a whole number.''
  \end{itemize}
\end{notebox}

\newpage
% ================================================================
\section*{Speaker 3 (Pritom) --- The World}

\begin{speakerbox}{Your goal in one sentence}
  Show that this is not a maths curiosity --- it is \stress{in nature, in your body,
  and in the software people use every day}.\\[4pt]
  {\small\color{DarkGray} Slides: Fractals in the Real World $\rightarrow$
  Conclusion (10 slides, 49 clicks) \quad$\cdot$\quad about 430 words}
\end{speakerbox}

\vspace{0.3cm}

We know what a fractal is, and we know how to measure one. So --- where is all of
this actually hiding?

\cue Everywhere in nature. Romanesco broccoli. \cue A snowflake. \cue An aloe.
\cue A pine cone. In every one of them, \stress{a small piece looks like the whole
thing}. \cue The same shape at every zoom.

\cue Here is that idea at its cleanest. Start with a square. Stand a
\stress{right triangle} on its top edge and grow a new square on each leg --- so
Pythagoras' theorem sits at every joint. \cue Three levels. \cue Seven.
\cue Twelve --- and a tree appears. \cue Nobody designed a single branch.

\cue Now watch what \stress{one number} does. At forty-five degrees the two
branches match exactly --- too perfect to be real. \cue Tilt that triangle to
fifty-five degrees \cue and they stop matching. \cue The tree \stress{leans}, the
way a real tree grows toward the light.

\cue And it is inside you. This is a cast of a lung --- the airways on one side,
the blood vessels on the other. \cue These are the vessels in a retina,
\cue and the same vessels zoomed in. \cue Your lungs split about
\stress{twenty-three times}, and that packs roughly \stress{seventy square metres}
of surface into about five litres of chest. \cue Doctors now measure the fractal
dimension of retinal vessels, because when that number drops it can be an
\stress{early sign of disease}.

\cue Now my favourite one. This fern is drawn by just \stress{four rules}, and
every dot is coloured by the rule that placed it. \cue Red is the stem.
\cue Green is the \stress{whole fern, shrunk and tilted} --- that one rule is
where all the detail comes from. \cue Gold is the left leaflet, \cue blue the
right.

\cue Four rules, six numbers each, plus four probabilities --- \stress{twenty-eight
numbers}. \cue And those twenty-eight numbers redraw the fern at \stress{any
resolution}. \cue Zoom in, and there is always more leaf. That is \stress{fractal
compression}: store the rule, not the pixels.

\cue Games use the same trick. One layer of bumps. \cue Three. \cue Seven --- and
a handful of numbers becomes a mountain range. \cue Push it to \stress{fifteen
layers} and you get ridges, valleys, snow lines and lakes. \cue Nobody drew any of
it.

\cue One honest catch. \stress{Real fractals stop.} A fern is a fern until you
reach a single cell. \cue Lungs stop at about twenty-three splits. \cue A tree
runs out at the last twig. \cue Only the mathematical ones go on forever.

\cue So, to bring it together. Part 1 gave us the surprise, \cue Part 2 gave us
the number, \cue Part 3 gave us the world. \cue Clouds are not spheres and
mountains are not cones --- but \stress{roughness has a number, and simple rules
make it}.

Thank you.

\begin{notebox}{Notes for delivery}
  \begin{itemize}[leftmargin=1.2em, itemsep=2pt, topsep=2pt]
    \item Your section is the one the audience will enjoy most. \stress{Let them
      look at the pictures} --- pause for a beat after each image appears.
    \item The lungs number (70 square metres inside 5 litres) always gets a
      reaction. Say it slowly.
    \item Your section has the most clicks of the three --- many are just the next
      picture appearing. \stress{Do not wait for each one}; keep talking and click
      as you go.
    \item On the four-rules fern, point at the \stress{green} part when you say
      ``the whole fern, shrunk''. That is the whole idea of the section.
    \item You close the whole talk. After ``Thank you'', \stress{stop and stand
      still}. Do not trail off or walk away mid-sentence.
    \item \stress{If you blank out}: ``Nature does not store details --- it stores
      a rule and repeats it. That is why fractals are everywhere.''
  \end{itemize}
\end{notebox}

\newpage
% ================================================================
\section*{If someone asks a question}

You are not expected to know everything. \stress{``That is a good question, I am
not sure''} is a completely acceptable answer and costs you nothing. But these
four come up often, so here are short answers.

\vspace{0.2cm}

\textbf{Is a fractal dimension a real dimension, like length and width?}\\
Not in the everyday sense. It is a measure of \stress{how much detail fills space
as you zoom in}. A shape with dimension 1.26 is rougher than a plain line but does
not fill an area.

\textbf{Are real things actually infinite fractals?}\\
No. Real objects are fractal only over a \stress{limited range of scales} ---
a coastline down to sand grains, lungs for about 23 branchings. Only the
mathematical ones are infinite. We covered this on the ``Real fractals stop'' slide.

\textbf{If the fern is only 28 numbers, why are image files so big?}\\
Because a photograph is not self-similar --- there is no small rule that redraws
your face. The fern works because it is \stress{built} from four copies of itself.
Real fractal compression hunts for self-similar \stress{blocks} inside an ordinary
image, which is far harder and much slower than JPEG.

\textbf{Is the terrain really made of 15 layers?}\\
Yes --- each layer is half the width and about half the height of the one before.
After about eight layers the detail is \stress{finer than a pixel}, so the last few
are there for completeness rather than for anything you can see.

\textbf{Is fractal compression actually used in practice?}\\
It was developed commercially in the 1990s and works beautifully on
self-similar images like the fern. For general photographs, JPEG and its successors
won out because they are much faster to encode. The \stress{idea} ---
describe the picture by the rule that makes it --- is what matters here.

\vspace{0.8cm}

\begin{notebox}{Before you walk in}
  Do one full run-through out loud, all three of you, with the slides. Not
  reading --- \stress{speaking}. It will feel awkward once and easy the second
  time. That single rehearsal is worth more than any amount of re-reading.
\end{notebox}

\end{document}
